\documentclass{article}
\IfFileExists{neurips_2025.sty}{\usepackage[preprint]{neurips_2025}}{\usepackage[margin=0.75in]{geometry}}
\usepackage{booktabs}
\usepackage{microtype}
\usepackage{amsmath}
\usepackage{hyperref}
\usepackage{xcolor}
\hypersetup{colorlinks=true, linkcolor=black, urlcolor=blue!55!black, citecolor=black}
\renewcommand{\thesection}{\arabic{section}.}

\title{A Nanochat-Style System for Fixed-Compute Protein Representation Learning}
\author{Lumin Science}

\begin{document}
\maketitle

\begin{abstract}
We describe a minimal, auditable system for training ESMC-class protein language
models under a fixed four-A100 compute budget. The system binds exact released
300M and 600M architectures to verified 70\%-identity cluster representatives,
two-stage source mixtures, hash-disjoint validation, benchmark decontamination,
FlashAttention execution, and frozen representation/contact evaluations. A
bounded diagnostic is used for routine iteration while the exact six-task
benchmark is reserved for releases. This paper is a living experiment report:
it distinguishes verified contracts from speedrun approximations and does not
fill unexecuted result cells.
\end{abstract}

\section{Motivation}
Large protein-language-model projects often intermingle corpus assembly, model
changes, optimizer tuning, and benchmark repair. We instead adopt one canonical
speedrun whose inputs, code, runtime, checkpoints, and evaluations are
content-addressed. The research objective is the best decision-grade protein
embedding model, not the lowest masked-token loss alone.

\section{Model and optimization}
Table~\ref{tab:models} fixes the two target tiers. Blocks use pre-normalization,
full-width Q/K normalization, non-interleaved rotary embeddings, rounded
$8/3$ SwiGLU, and residual updates divided by $\sqrt{N/36}$, following the
released Biohub implementation. The paper text instead states $\sqrt{N}$; we
record rather than conceal this discrepancy.

\begin{table}[h]
\centering
\caption{Frozen architecture contracts.}
\label{tab:models}
\begin{tabular}{lrrrrr}
\toprule
Tier & Layers & Width & Heads & FFN & Parameters \\
\midrule
ESMC 300M & 30 & 960 & 15 & 2560 & 332,997,184 \\
ESMC 600M & 36 & 1152 & 18 & 3072 & 575,036,992 \\
\bottomrule
\end{tabular}
\end{table}

Training uses 15\% mask-all corruption, sequence-mean cross entropy, AdamW with
$\beta=(0.9,0.95)$, selective decay, and unit gradient clipping. The original
warmup--stable--decay trajectory is compressed into 1,800 measured seconds:
1,200 seconds at context 512 and 600 seconds at context 2,048. Evaluation and
checkpoint I/O are excluded from this training clock. Step duration is
max-reduced across DDP workers so every rank makes the same context-stage and
stopping decision.

All Python dependencies are managed by uv with a committed lockfile. The
qualified Linux environment fixes Python 3.11 and Torch 2.13 with CUDA 12.6.
Prefix-padded batches are compacted with cumulative sequence lengths and passed
to Torch's in-tree variable-length FlashAttention operator, avoiding both
quadratic padding work and long-sequence-only sampling bias.

The operational sweep increased the per-GPU microbatch from 8/1 to 64/16 for
the 512/2,048-token stages. In a 60-second qualification this increased model
tokens from 2.096M to 4.521M (2.16$\times$) and sequences from 8,128 to 17,408.
Peak observed throughput was 101.8k non-pad model tokens/s (16.3\% 6N MFU) in Stage 1
and 46.1k/s (7.4\%) in Stage 2, relative to the A100 dense BF16 peak of 312
TFLOP/s. Dynamic compilation was rejected after an AOTAutograd metadata failure
before step one; eager execution is therefore the qualified canonical path.

\section{Data}
The verified transfer contains 92.2M UniRef, 348.1M MGnify, and 324.9M OMG/IMG
representatives. Because member sequences were not transferred, sampling is
uniform over representatives rather than the paper's cluster-then-member
procedure. Stage mixtures are 36/11/54\% and 63/6/31\% respectively. Exact
sequence hashes in P-CORE (108,215 sequences) and the contact manifest (20,795
chains) are excluded; homology-level decontamination remains a release blocker.

\section{Evaluation}
The routine profile reports held-out sequence-mean MLM NLL, long-range contact
precision at $L$, and three protein-level representation task metrics: remote
homology, human PPI, and FLIP2 fitness. The latter are exact task metrics from
P-CORE v0.2 but omit bootstrap and do not form a P-CORE aggregate. Each probe
runs in an isolated process with a ten-minute limit, and timed-out tasks remain
explicitly missing. P@L uses a logistic probe over symmetrized attention planes
from every head/layer, 20 frozen probe structures, an 8\AA{} C$\beta$ threshold
(C$\alpha$ for glycine), and sequence separation at least 24.

Full P-CORE v0.2 is a release evaluation, not a 30-minute-training gate. The
original serial evaluator embedded residues for all 108,215 sequences, produced
a 122~GB cache, and spent more than one hour in four secondary-structure LBFGS
fits. The optimized exact path stores protein means for all tasks and residue
embeddings only for the 11,411 secondary-structure sequences. Release probes
remain exact and should execute as restartable task-parallel jobs.

The Enzyme Commission task is quarantined for model selection. Released 300M
and 600M models score approximately 71 task skill, whereas 6B scores 1.607 even
though an integrity audit found complete, finite, non-degenerate embeddings and
disjoint splits. The release profile executes and reports EC but does not allow
it to decide a model until the failure is reproduced and repaired.

\section{Current reference and speedrun results}
The current released-model reference is P-CORE 45.550 for ESMC-300M and 45.690
for ESMC-600M; the paired 95\% interval for the difference includes zero.
On a frozen 1,024-chain diagnostic, P@L is 0.5340 and 0.5778 respectively.

\begin{table}[h]
\centering
\caption{Canonical local speedrun. Training throughput uses non-padding model tokens.}
\begin{tabular}{lrrrr}
\toprule
Checkpoint & Train seconds & Tokens/s & Held-out NLL & P@L \\
\midrule
Stage 1 & 1200.317 & 93,099 & 2.7289 & 0.0523 \\
Final & 1800.024 & 75,196 & 2.7544 & 0.0619 \\
\bottomrule
\end{tabular}
\end{table}

The held-out NLL estimate uses 32 validation sequences and the contact diagnostic
uses 32 predeclared chains, so neither supports a small-difference claim. Stage~2
improves diagnostic P@L while its NLL is slightly worse; this divergence is a
reason to retain downstream probes rather than tune on MLM loss alone.

\begin{table}[h]
\centering
\caption{Bounded representation diagnostic; no aggregate is defined.}
\begin{tabular}{lrrr}
\toprule
Checkpoint & Remote homology BA & Human PPI AP & FLIP2 $\rho$ \\
\midrule
Stage 1 & 0.0682 & 0.8519 & 0.3162 \\
Final & 0.0724 & 0.8574 & 0.3150 \\
\bottomrule
\end{tabular}
\end{table}

The final checkpoint is directionally better on remote homology, PPI, and P@L,
but slightly worse on FLIP2 and MLM NLL. These single-run diagnostics therefore
do not support a general quality claim. With cached MLM/contact components, both
checkpoints completed concurrently in 8.81 minutes with all three tasks present;
the measured components imply roughly 16 minutes for a cold routine evaluation.
The prior exact full-profile attempt exceeded one hour in secondary structure
alone, which motivated the profile separation.

\section{Planned experiments}
We first select operational settings (microbatch, accumulation, activation
checkpointing, compilation) by throughput and loss parity. We then calibrate
the undisclosed proxy learning rate/decay and the residual-scaling discrepancy
at small scale before multi-seed 300M confirmation. Only after evaluation repair
and homology decontamination do we promote source-mixture, context, or 600M
comparisons to scientific claims.

\section{Limitations}
A 30-minute run cannot approximate the original 1.5M-step optimization horizon.
Representative-only sampling differs from the released data process; exact
hash exclusion does not prevent homologous contamination; one P-CORE task is
untrusted; the routine three-task vector is not a P-CORE score; and the contact
speedrun subset is diagnostic rather than the full 20,775-chain evaluation.
These are explicit promotion gates, not footnotes.

\section*{Primary sources}
Code and architectural contracts are frozen against Biohub ESM commit
\href{https://github.com/Biohub/esm/commit/9b9078cadf4e08bd02e9715ba99313fc5127379a}{9b9078c}.
The training and contact protocols follow the 2026 ESMC preprint,
\href{https://doi.org/10.64898/2026.06.03.729735}{doi:10.64898/2026.06.03.729735}.
Repository organization is inspired by
\href{https://github.com/karpathy/nanochat}{karpathy/nanochat}.

\end{document}
